\begin{figure}[tbp]
\centering
\begin{tikzpicture}[x=1cm,y=1cm]
\node[note,anchor=west] at (-1.0,0.9) {普通嵌套：以一个需要 L1 处理的退出事件为例};
\node[box,text width=1.5cm] (a) at (0,0) {L2};
\node[box,text width=1.5cm] (b) at (3,0) {L0};
\node[amberbox,text width=1.5cm] (c) at (6,0) {L1\\处理事件};
\node[box,text width=1.5cm] (d) at (9,0) {L0};
\node[box,text width=1.5cm] (e) at (12,0) {L2};
\draw[wire] (a)--node[above,font=\footnotesize]{退出}(b);
\draw[wire] (b)--node[above,font=\footnotesize]{反射}(c);
\draw[wire] (c)--node[above,font=\footnotesize]{请求恢复}(d);
\draw[wire] (d)--node[above,font=\footnotesize]{进入}(e);
\node[note,anchor=west] at (-1.0,-1.15) {HyperTurtle：已注册的处理逻辑在 L0 执行};
\node[greenbox,text width=1.5cm] (f) at (0,-2.1) {L2};
\node[greenbox,text width=5.0cm] (g) at (6,-2.1) {L0：调用 L1 的 eBPF hyperupcall\\共享 maps、受限 helper、状态同步};
\node[greenbox,text width=1.5cm] (h) at (12,-2.1) {L2};
\draw[wire] (f)--node[above,font=\footnotesize]{退出}(g);
\draw[wire] (g)--node[above,font=\footnotesize]{成功后恢复}(h);
\node[amberbox,text width=3.1cm] (fallback) at (6,-4.0) {L1 原有处理路径\\保留正确性与复杂功能};
\draw[dwire] (g)--node[right,font=\footnotesize]{资源不足 / 竞争 / 不支持时回退}(fallback);
\end{tikzpicture}
\caption[HyperTurtle 减少跨层往返]{HyperTurtle 减少跨层往返的方式。图显示通用热路径，具体 EPT 缺页还可能包含重执行后再次退出等步骤；成功下放只减少特定处理所需的切换，不取消 L1。\cite{source8}}
\label{fig:hyperupcall}
\end{figure}
